\documentclass[UTF8,zihao=-4]{ctexart}
\usepackage[a4paper,top=2.5cm,bottom=2.5cm,left=2.5cm,right=2.5cm]{geometry}
\usepackage{booktabs,longtable,array}
\usepackage{setspace}
\usepackage{hyperref}
\hypersetup{hidelinks,pdfauthor={},pdfsubject={CUMCM 2026 AI 工具使用详情}}
\ctexset{
  section={format=\zihao{3}\bfseries,beforeskip=1.8ex,afterskip=0.9ex}
}
\setstretch{1.25}
\setlength{\parindent}{2em}
\setlength{\emergencystretch}{3em}

\title{AI工具使用详情}
\author{}
\date{}

\begin{document}
\maketitle

% 本文件属于支撑材料，也不得出现参赛者、学校和赛区身份信息。

\section{所用工具}

\begin{longtable}{p{0.28\textwidth}p{0.28\textwidth}p{0.34\textwidth}}
\toprule
工具名称 & 版本或型号 & 使用日期/范围 \\
\midrule
请填写 & 请填写 & 请填写 \\
\bottomrule
\end{longtable}

\section{具体使用目的和环节}

请填写。逐项说明用于语言润色、资料整理、代码调试、建模讨论或其他环节的具体情况。

\section{主要提示方式与使用过程}

请填写。概括主要提示策略、输入信息与交互过程；必要时附典型交互示例，但不得泄露身份信息。

\section{采纳、人工修改与核验}

请填写。逐项记录采纳了哪些输出、如何修改、使用何种独立计算或资料进行核验，以及核验结果。核心建模与分析必须由参赛队主导。

\subsection*{2026-09-10：研究路线文档整理记录}

本次使用 Codex 辅助建模讨论与项目文档整理，具体会话型号待参赛队核对补充。输入为参赛队提供的第一性原理思考、项目工作规范、题面和参数配置；任务为统一四问的研究主线。

输出为 \texttt{output/research-roadmap.md}，包含统一物理与信息模型、费用定义、四问递进路线、实验设计和实施标准。整理过程中核对了原始题面 PDF，区分了题面事实与建模假设，并补充了充放电互斥、非对称费用推导条件、合同结算解释、因果回测及终端价值的限制。未进行调度求解，未生成模型数值结论，也未开展外部文献检索。

本记录由工具按实际操作生成；研究路线是否采纳、具体人工修改、公式独立核验及核验结果均待参赛队填写。本记录不代表已经完成人工审查。

\subsection*{2026-09-10：实验前论文理论稿撰写记录}

根据参赛队“暂不开展实验，先完成问题分析等可写部分”的要求，使用 Codex 将研究路线整理为数学建模竞赛论文初稿。新增问题重述、分问分析、假设、符号说明、统一物理与信息模型、四问理论表达、评价方法及适用范围，调整摘要、标题和题面引用。未运行预测、调度、回测或敏感性实验，结果、算法参数及经验性结论保留待补充标记。

本次技术检查包括 LaTeX 编译、交叉引用和 PDF 排版检查，不等同于数学模型的独立数值验证。正文采纳情况、公式核验、假设确认及人工修改仍待参赛队记录；本次没有新增外部文献检索。

\subsection*{2026-09-11：第一问求解、校验与可视化记录}

根据参赛队完成第一问并保留结果供后续精进的要求，使用 Codex 辅助实现共享电池参数、确定性整数优化、费用优先的分层求解、线性松弛与对偶校验、独立能量审计、官方工作簿写入、基线及参数敏感性、版本归档、论文图表和离线交互图。原始附件及模板按只读处理，参赛队已明确选择模板时段起点及午夜日循环的解释，另一时间口径另存对照。

本次计算得到主模型最低购电费约35126.85元，并将验证后的结果写入论文。自动核验包括可手算小系统、物理与费用重算、CSV及Excel回读、模板格式、线性松弛下界和输入及产物哈希；这些属于程序校验，不代替参赛队的独立人工审查。求解环境存在的库导入警告记录于运行manifest，具体版本、参数、源码和数值容差随结果保存。实现中参考了SciPy官方求解接口文档，已在论文引用。

本次AI输出的具体采纳、人工修改、模型假设确认及独立复核情况，仍由参赛队据实际情况填写。

\clearpage
\subsection*{2026-09-11：第二问因果回放、比较与论文更新记录}

根据参赛队继续完成第二问的要求，使用Codex辅助实现仅依赖历史的负荷与光伏预测、样本外残差分位数、固定日前合同、三类因果储能控制器、跨日连续回放、物理与费用独立校验、模板扩展、版本归档及可视化。参赛队明确同意将第一问的日循环起点解释延续到第二问。

开发中先比较滚动控制与简单控制表现，再将全部控制器纳入一月验证排序；最终72个候选组合仅按一月分数选定，二月至十二月作为因果历史回放，不冒称未接触的独立测试集。主方案采用趋势预测、80\%分位数与贪心缓冲，实际总费用约14289973.02元。平均费用优势、尾部风险不全面占优及参考SOC跟踪的小幅差异均如实保留。

自动验证包括小型手算案例、未来实测扰动下的前缀不变性、全部运行轨迹的能量及费用重算、求解原始--对偶核验、Excel格式与数值回读，以及源码和产物指纹。程序验证不替代参赛队人工审查。已参考并引用滚动预测验证教材与SciPy官方文档；PDF技能用于原题核对、论文编译后的页面渲染及排版检查。环境导入警告仍记录在运行材料中。

参数和模型是否采纳、核心判断的独立复核及具体人工修改情况，仍需参赛队据实补充。

\subsection*{2026-09-11：科研图形统一与图文同步更新}

根据参赛队要求，使用Codex辅助建立Nature及NeurIPS启发的科研绘图规范。查阅Nature官方图形指南和NeurIPS 2023官方格式示例，保存出处于\texttt{code/vis/STYLE\_GUIDE.md}；仅参考清晰性、黑白可读和矢量输出原则，不将其当作本竞赛投稿格式。

建立集中样式配置、共同绘图接口、最小字号与画布边界检查、PDF/SVG/PNG三格式输出及独立视觉归档。重绘问题一、二现有图形，并从已保存的比较统计生成费用--尾部风险及配对节费区间图；同步补充论文图注、收益边界分析、源程序附录和离线图形索引。只改变呈现，不重算优化、不平滑真实轨迹、不修改原数值归档或工作簿；对保护文件核对前后哈希。

自动样式测试与PDF页面渲染用于检查实现和排版，不等于人工审查。图形是否采纳、最终审美判断、公式及统计解释的独立复核，仍由参赛队填写确认。

\subsection*{2026-09-11：问题二跨日边界B方案修订}

参赛队审阅日末储能处理后明确要求采用48小时方案。使用Codex辅助将预测扩展为共享历史截止点的直接多提前期预测，按目标日已结束原则校准相应提前期残差；扩展已有LP并仅提交前144时段合同，保留次日辅助轨迹、实际SOC连续和共同残值记账。取消主模型旧终端系数网格，以24组一月候选重新选型；远端零价值作为明确近似，以72小时、固定远端残值及24小时替代成本作完整回放对照。

核对发现，旧版终端系数0.8与1.0并列，不能宣称前者独立占优；完整10分钟电价的最低点在05:40，不能用整点价推导全天阈值。新版总费用14304319.83元，比旧版增加14346.81元，晚间补购显著减少；未将机制改善误写成费用必然降低。新版2月1日日末电量高于6000 kWh，与旧版当日消耗初始电量的情况不同。两版实际结果及原始归档均保留。

同步更新论文、表格、Excel、统一风格图及运行说明。验证包括多提前期未来扰动、成熟残差、首日合同截取、跨日SOC连续、年末视窗、全部实际与辅助轨迹的独立审计，以及预测因果重建。PDF技能用于编译后的页面检查。该修订由已经查看的年度表现触发，不宣称未接触测试集。此记录为工具执行记录；核心假设、采纳情况和独立人工复核仍需参赛队填写。

\subsection*{2026-09-12：问题三发布预报接入与有偿调整}

根据参赛队要求，使用Claude Code辅助完成第三问。核对附件三结构后，用四种对齐假设与附件二实测比较，确定“预报$h$小时”对应发布后$h$小时起始的10分钟时段，以发布前刚结束时段实测为首锚点插值；融合权重0.5是在查看全年附件三、本地趋势预测与等权融合的光伏误差诊断之后、在任何费用回放之前确定的，此后未按费用结果调整，故“先验”指相对费用回放而言，不是相对预测误差诊断而言。结算主口径改为字面读法的退款口径（实际合同按原价、偏离另付50\%电价），原路线中的不退款口径改为对照；该改动由工具提出并在记录中说明理由，采纳与否由参赛队确认。

在第二问48小时规划与贪心执行器上增加有偿调整线性规划，调整带$[0.7,0.9]$由“1减边际获取价与紧急价之比”导出，不作全年调参。主策略先验固定为三时点全部调整；时点子集、预报冻结、仅附件3、无附件3、结算口径与对称规则共14种策略全年回放，一月15--31日仅作规则检验。主策略现金费用13712887.43元，较仅计划节省449572.89元，较第二问节省591432.40元；发现节省主要来自调整权限与状态反馈，最新光伏预报增量在统计上不显著，如实写入论文。

自动验证包括对齐与因果锚点测试、调整带与结算规则手算案例、发布后观测改写下的前缀不变性、第二问信息集下的逐时复现、全部轨迹的物理与三项费用重算、调整记录与因果重建核验、工作簿回读及产物指纹。PDF技能用于编译后的页面检查。此记录为工具执行记录；核心判断与独立人工复核仍需参赛队填写。

\subsection*{2026-09-12：问题四实时电价模型}

根据参赛队要求，使用Claude Code辅助完成第四问。核对附件四后确认价格均值与附件一相同、日内平均形态与附件一完全相关、日级水平随星期与净负荷变动。在多种因果预测器中，先以一月8--31日的价格预测误差检验比较，再固定两周同日均值与净负荷回归的等权融合；该组合相对二月至十二月的费用回放是先验的，相对一月误差诊断不是，此后未按费用结果调整。日内修正取已观测价格的平均偏差加性修正。将第三问日回放函数扩展为接受价格模型：规划用预测价，全部账单按交付时段附件四价格结算，价格模型为常数时逐时复现第三问。

以第二问与第三问信息集分别全年回放，另设附件一固定轮廓、当日价格已知、无日内修正、仅计划对照，共8种策略。4-2主策略现金费用14986891.98元，4-3主策略14383873.10元；发现相同决策仅因实时价结算即多付约5.4\%，价格预测收回约0.5--0.6\%，日内价格修正无显著作用；当时把4-3相对4-2的全部差额归于“合同调整机制”，第一轮审计（R10）后已拆分为0:00融合外部预报与日内有偿调整两部分（见下文），均如实写入论文。

自动验证包括价格模型的因果性与观测边界测试、附件一常数模型对第三问的兼容复现、实时价结算恒等式、全部轨迹的物理与费用重算、规划价格的因果重建与未来价格改写检验、两份工作簿回读及产物指纹。第三问归档因回放函数扩展而重跑，数值完全一致。此记录为工具执行记录；核心判断与独立人工复核仍需参赛队填写。

\section{人工责任确认}

待参赛队完成人工审查后填写确认情况，并对参赛作品的原创性、真实性和准确性负责。本模板及上述自动记录不构成人工确认。

\subsection*{2026-09-11：正文图形全面重做与图形自动审计记录}

根据参赛队“全面优化正文绘图、可补充新图、兼顾美观与新颖并适配问题”的要求，使用Claude Code（模型Claude Fable 5.1）及其科研图形技能，对第一至四问正文图重新设计并新增一张四问信息—决策—结算时间线示意图。设计依据为“一图回答一个结果层面的问题、各子图承担不同推断角色”的原则：第一问以充放电底色贯通四个面板，并以电价与LP边际成本之差显示储能替代购电的时段；第二问以指定日期中紧急购电最多的一天展示情景带、风险轮廓与合同，同图给出费用构成、配对节费区间与尾部风险；第三问以费用阶梯分解外部预报、调整权限与最新预报的价值，并补充时点边际价值、调整足迹与预报误差；第四问以星期与时段热图、日均价与净负荷散点展示价格结构，以双费用阶梯与括线分解结算、预测、信息与机制效应。示意图不含实际数值，图注已注明。

所有图形仍从已校验的数值归档重绘，构建流程逐文件核对归档哈希，本次未修改任何模型、参数或结果；工作簿与表格数字不变。为控制图形质量，在绘图流程中加入渲染后面板对齐审计（共享边、宽高与间距容差1.5 pt）、PDF最小字号审计与PDF文字—线条、文字—文字交叠审计，并逐面板目视检查；审计中排除了PDF阅读库对纹理图案单元的误报，第一问图中图例与底色的交叠为有意保留并已核对。旧的费用比较图与附录风险权衡图由新图取代并从发布目录移除；图注、正文交叉引用及风格规范同步更新。

图形是否采纳、图注表述及与正文数字的一致性核验，仍由参赛队据实复核填写。

\subsection*{2026-09-11：第一轮建模方法与求解审计}

根据参赛队要求，使用Codex辅助审阅四问的物理模型、信息边界、风险代理、结算、求解与验证流程，形成\texttt{doc/Review-1st.md}。本轮不修改模型、参数、论文结论或官方结果。重新运行57项测试及四问独立结果核验，全部通过；另以输入扰动确认第三问日内发布预报的跨日部分未参与当前决策，以小型算例检查随机价格分位数、日前调整期权和调整带内价值函数的适用范围。

审计区分已确认问题与待实验验证的改进：当前物理及费用核验通过，不代表风险轮廓代理等价于原随机控制最优模型；价格预知的启发式对照也不构成严格最优费用界限。参考SciPy官方文档、滚动预测验证教材、滚动控制及CVaR方法的一手资料，链接保存在审计报告；PDF技能用于原题页面核对。本轮新增的诊断结果不是新策略的年度收益，不改写既有数值归档。建议是否采纳及核心结论的独立人工复核仍由参赛队确认。

\subsection*{2026-09-12：第一轮审计的处理（迭代0与迭代1）}

根据参赛队要求，使用Claude Code（模型Claude Fable 5.1）逐项核实\texttt{doc/Review-1st.md}的R01--R13并实施迭代0（纠正论证）与迭代1（补全信息链），处理记录见\texttt{output/review-1st-response.md}。核实结论：R02、R07、R13为已确认且尚未修复；R01、R04、R05、R06为已确认的论证问题，以小型反例测试固定并改写正文，算法本身未升级；R03、R08、R09、R10、R11、R12为合理近似或范围扩展事项，本轮只做说明、归因拆分与记录补全，真实费用目标、调整期权与带内价值函数优化列为下一阶段。未发现审计判断需要修正之处；补充确认价格预知对照在4-2与4-3各有132/334天逐日费用高于因果预测策略。

代码改动：第三问在每次调整时把当前发布版本的次日覆盖部分接入次日辅助段，覆盖之外用0:00本地两日提前期预测补齐，残差按同一构造、只用成熟目标重建；日内负荷偏差修正的校准残差重建各历史日期在同一发布时刻本可执行的修正；两个开关均关闭时逐时复现2026年9月11日冻结归档（最大差$10^{-12}$量级）。0:00计划的288个时段与每次调整的完整辅助段被保存，并通过物理、承诺标记与归档配置下的因果重建核验；校验器改为按归档配置重建并增加未发布版本与未来实测的毒化检查。第四问增设冻结旧版信息链对照，“当日价格已知”改称启发式对照，4-3相对4-2的差拆为0:00融合预报与日内有偿调整两部分。新增18项测试（信息链13项、审计反例5项），全部75项测试通过。

结果：第三问主策略现金费用由13712887.43元变为13701007.35元，减少11880.08元（0.09\%），日均差95\%区间含零；差额几乎全部来自同构残差，跨日新预报单独作用为几元量级且方向为略增（仅跨日相对旧版$+3.84$元，两项同用相对仅同构$+7.21$元）；紧急购电由37589.90增至39662.16 kWh，CVaR95由6557.78增至6812.62元。第四问4-3由14383873.10元变为14373781.70元，4-2不变。新版本被采纳为正式发布的依据是事先约定的“信息集声明与实现一致”，不是费用更低；费用未显著下降、尾部风险略升、各版本次日覆盖段光伏误差高于本地预测、单一分位数对照在本年度更便宜等负面或不利结果均如实写入论文。工具同步更新了研究路线、Agent规范、README、论文四问相关章节、图注与本记录；PDF技能用于编译后的页面检查。此记录为工具执行记录；核心判断、方案采纳与独立人工复核仍需参赛队确认。


\section{图一、图六重绘与全文润色（2026年9月12日）}

用户要求重新设计图一为论文主图、美化图六，并润色全文、将正文控制在30页以内。Codex使用PDF技能检查论文与矢量图源代码，采用Matplotlib重绘：图一按微网物理能量流、因果调度闭环与四问递进组织；图六将绝对费用、相邻变化、紧急购电和4-3相对4-2的收益分解分面呈现。此次未调用图像生成模型，图中文字与数值均由矢量代码输出，费用与区间读取已归档数据，不修改模型计算结果。

文字修改集中于合并重复背景、精简统一模型与各问说明、统一调整变量符号、澄清费用比较的百分比分母，并把详细因果核验移入附录。保留指定日期结果、非显著或不利对照、信息可用边界及代理模型局限。重新编译论文与本记录，检查正文页数、交叉引用、图内字号、裁切与重叠；视觉归档保留数值与工作簿哈希。该记录描述工具完成的编辑与核验，不表示参赛队已完成人工审阅。


\section{图一改用图像生成工具绘制（2026年9月12日）}

根据用户进一步要求，以内置GPT图像生成工具重新绘制图一，采用电网、光伏板、储能柜与住宅的技术插画，分层展示微网能量流、信息约束下的调度闭环和四问递进。工具接口不提供模型版本选择，也未返回可核实的具体版本，故不能将本次调用标记为GPT Image 2.5。原图生成后，再次调用图像工具修正储能充、放电箭头和“光伏预测”标签；未使用程序改画生成图。

最终PNG原生尺寸为1448乘1086像素，在160 mm图宽下约230 dpi。配套PDF仅封装原始位图，不是矢量图，也未将插值放大标为高分辨率原图。原矢量版保留作可编辑备选。完整提示词、初稿、定稿与版本信息保存于\path{output/imagegen/}，论文引用\path{output/figures/framework_imagegen.png}。核对项目包括电网至微网的单向取电、储能双向能量流、当前可用信息、日内调整时点、设备边界及第四问的两类合同权限；重新编译后检查插图可读性与30页主文限制。该核验为工具检查，未代替参赛队人工复核。

\section{附录精简与复现索引（2026年9月12日）}

根据用户明确要求，将论文完整版本从逐文件打印源码改为精简附录。保留补充图表、指定日期结果、信息链对照及独立核验，删除完整代码清单的排版引用，补充运行环境、复现命令、核心文件与结果索引，并直接截取四问的关键函数代码，注明原文件行号和依赖边界。完整模型、数据处理、绘图及测试源码继续保存在项目目录，本次不改动模型算法或数值结果。重新编译后核对全文约50页、附录前30页内容一致、公式与图表交叉引用及页面可读性。此次精简由用户指定，不表示原完整源码版已被删除，也不表示已进行人工复核。

\section{写作结构重组与论证润色（2026年9月12日）}

根据用户对问题分析偏简、后续章节冗杂的反馈，Codex重新组织论文论证。第一章扩充研究任务、储能时序耦合、误差成本不对称、信息与交易权限的关系，并将四问的方法选择与检验思路对应说明。第四章集中共同物理、信息及经济边界；各问精简重复的通用目标，保留实际求解模型、关键推导和指定结果。第三问围绕收益归因与时点消融组织证据，第四问区分结算效应、因果价格预测和有偿调整收益，第九章集中核验、对照解释与适用边界，结论改为机制归纳。

本轮为结构和文字编辑，未改变模型源码、实验参数、数值归档、图形及结果工作簿；保留未显著节费、尾部风险升高、单一分位数对照更便宜等不利结果，不将代理求解写成随机控制全局最优。使用PDF技能重新编译并渲染检查，更新主文阅读版与完整论文，核对主文不超过30页、全文约50页及公式和图表引用。旧稿和核验记录保存在独立写作版本目录。此记录不代表竞赛获奖等级判断，也不代替参赛队人工审阅。

\section{第一章扩充、主图定版与全文语言润色（2026年9月12日）}

根据用户要求，Claude Code重写第一章为“问题背景与重述、问题分析、总体思路”三节：补入储能参数与四问的逐条重述，将主图（图一，图像生成版）置于第一章末，并以强制换页使第一章恰好占据第2至3页。第九章更名为“模型检验、评价与推广”，新增“模型的优点”小节并在局限段补充推广方向；问题三新增对“是否需要引入其他时刻预报”的直接回答段。全文统一“问题一至四”的称谓，润色摘要、统一模型、四问与结论的表述，改善段首主题句与章节衔接。

本轮只改动论文文字、章节结构与图一版面（图宽调至0.95倍版心、图注压缩），未改变模型源码、实验参数、数值归档、图形数据与结果工作簿；所有数值仍由\path{output/tables/}宏引入。重新编译后核对：正文止于第28页、参考文献第29页、全文49页，无溢出框与未定义引用，\texttt{make check}通过。该记录描述工具完成的编辑与核验，不代替参赛队人工审阅。

\end{document}
